%% kapa: exam=1 round=30 year=2019 period=2 form=A subject=appraisal_law subject_ko=감정평가관계법규 seq=16 number=16
%% source: Q-Net 감정평가사 기출문제 (cst00309) artlSeq=5205205 file=extracted/1008/5205205/2019년 제30회 감정평가사 1차시험 문제지/2019년 제30회 감정평가사 1차시험 문제지/2019년도 제30회 감정평가사 1차 2교시 A형.hwp
%% conversion: kordoc 4.24.51 -> tools/gen1.py (auto); review=proofread-sonnet; flags=none
\begin{question}{16}
감정평가 및 감정평가사에 관한 법령상 감정평가업자에 관한 설명으로 옳은 것은?
\begin{choices}
  \item 감정평가업자가 토지를 감정평가하는 경우에는 그 토지와 이용가치가 비슷하다고 인정되는 토지의 적정한 실거래가를 기준으로 하여야 한다.
  \item 감정평가업자는 감정평가서의 원본을 발급일부터 2년 동안 보존하여야 한다.
  \item 감정평가업자는 토지의 매매업을 직접 할 수 있다.
  \item 감정평가업자의 업무수행에 따른 수수료의 요율은 국토교통부장관이 감정평가관리ㆍ징계위원회의 심의를 거치지 아니하고 결정할 수 있다.
  \item 감정평가업자가 감정평가서 관련 서류를 보관하는 기간과 국토교통부장관이 감정평가업자의 해산이나 폐업시 제출받은 감정평가서 관련 서류를 보관하는 기간은 동일하다.
\end{choices}
\end{question}
\answer{5}
